\documentclass[border=5mm]{standalone}
% ==============================================================================
% SETUP.TEX - CONFIGURACIÓN GLOBAL DEL PROYECTO
% ==============================================================================
% Este archivo centraliza todos los paquetes y estilos.
% Debe incluirse en el preámbulo del libro principal y de los archivos standalone.
% \PassOptionsToPackage{gray}{xcolor}
% --- 1. CODIFICACIÓN E IDIOMA ---
% fix-cm habilita las fuentes Computer Modern en tamaños arbitrarios (por debajo
% de 5 pt, entre otros). Debe cargarse antes que fontenc.
\usepackage{fix-cm}
\usepackage[utf8]{inputenc}
\usepackage[T1]{fontenc}
\usepackage[spanish, es-tabla]{babel}  
\usepackage{csquotes} % Recomendado para babel y citas

\usepackage{import}
% mode=image: Usa el PDF pre-compilado, nunca intenta recompilar.
% Debes compilar cada figura manualmente antes de compilar el libro.
\usepackage[mode=image, subpreambles=false]{standalone}

% --- 2. MATEMÁTICAS Y CIENCIAS ---
\usepackage{amsmath, amssymb, amsfonts, mathtools}
\usepackage{enumitem}

% LaTeX trae \DeclareMathSizes{5}{5}{5}{5}, así que a 5 pt (\tiny) los subíndices
% salen del mismo tamaño que la base: en las etiquetas de figuras el M de
% $\omega_M$ queda desproporcionado. Se restablece la proporción habitual.
\DeclareMathSizes{5}{5}{3.7}{3}

% --- 3. DISEÑO Y GEOMETRÍA --- 
% Unificamos los márgenes para todo el documento
\usepackage[left=2.5cm,right=2.5cm,top=3cm,bottom=3cm]{geometry}
\makeatletter
\ifstandalone % Evita que geometry dañe el recorte de las figuras bajándolas 3cm
  \@ifclasswith{standalone}{crop=false}{
    % Es un capítulo/sección: Dejamos que geometry actúe.
  }{
    % Es una figura: Neutralizamos geometry para que no las recorte mal.
    \geometry{pass}
  }
\fi
\makeatother
\usepackage{parskip} % Espaciado entre párrafos sin sangría
\usepackage{float}   % Para usar [H] en figura

% Ajustes de flotantes para evitar espacios verticales exagerados
\renewcommand{\topfraction}{0.95}      % Permite que las figuras ocupen hasta el 95% superior
\renewcommand{\bottomfraction}{0.95}   % Permite que las figuras ocupen hasta el 95% inferior
\renewcommand{\textfraction}{0.05}     % Permite hojas con tan solo un 5% de texto
\renewcommand{\floatpagefraction}{0.8} % Requiere que una página de solo figuras esté 80% llena
\setcounter{topnumber}{4}
\setcounter{bottomnumber}{4}
\setcounter{totalnumber}{10}

% --- 4. GRÁFICOS Y TABLAS ---
\usepackage{graphicx}
\usepackage{booktabs} % Tablas profesionales
\usepackage{caption}  % Control de captions y \captionof
\usepackage{tikz}
\usetikzlibrary{arrows.meta, calc, angles, quotes, babel, backgrounds}
\usepackage{pgfplots}
\usepgfplotslibrary{groupplots}
\usepgfplotslibrary{external}
\usepgfplotslibrary{patchplots}
\pgfplotsset{compat=1.18}
\usepackage[american, straightvoltages]{circuitikz}

% --- 5. COLORES Y CAJAS ---

\usepackage[table]{xcolor}
\usepackage{tcolorbox}

% Definición de colores corporativos/estilo
\definecolor{utnblue}{RGB}{0, 51, 102}
\definecolor{codegreen}{rgb}{0,0.6,0}
\definecolor{codegray}{rgb}{0.5,0.5,0.5}
\definecolor{codepurple}{rgb}{0.58,0,0.82}
\definecolor{backcolour}{rgb}{0.96,0.96,0.96}

% --- 6. ENTORNOS PERSONALIZADOS (CAJAS) ---

% Caja "Resultado" (Azul Corporativo) — Definiciones, fórmulas clave, resultados importantes.
% Uso: \begin{resultado}[Fórmula de Euler] ... \end{resultado}
\newtcolorbox{resultado}[1][]{
    colback=utnblue!5!white,
    colframe=utnblue,
    title=\textbf{#1},
    fonttitle=\bfseries,
    sharp corners,
    boxrule=0.5mm
}

% Caja "Nota" (Gris) — Advertencias, aplicaciones, contexto secundario.
% Uso: \begin{nota}[Cuidado con la calculadora] ... \end{nota}
% Uso: \begin{nota} ... \end{nota}  → título por defecto "Nota"
\newtcolorbox{nota}[1][Nota]{
    colback=codegray!5!white,
    colframe=codegray,
    title=\textbf{#1},
    fonttitle=\bfseries,
    sharp corners,
    boxrule=0.5mm
}

% Caja inline para resaltar un valor y enlazarlo con su otra aparición en una tabla
% o en una cuenta: mismo color, mismo valor.
% Uso: \marcavalor{$4$}  o  \marcavalor[codegreen]{$4$}
\newtcbox{\marcavalor}[1][utnblue]{
    on line,
    boxsep=1pt, left=2pt, right=2pt, top=1pt, bottom=1pt,
    colback=#1!10!white,
    colframe=#1,
    boxrule=0.4pt,
    arc=2pt
}

% --- 7. CÓDIGO FUENTE (LISTINGS) ---
\usepackage{listings}

\lstdefinestyle{miestilo}{
    backgroundcolor=\color{backcolour},   
    commentstyle=\color{codegreen},
    keywordstyle=\color{magenta},
    numberstyle=\tiny\color{codegray},
    stringstyle=\color{codepurple},
    basicstyle=\ttfamily\footnotesize,
    breakatwhitespace=false,         
    breaklines=true,                 
    captionpos=b,                    
    keepspaces=true,                 
    numbers=left,                    
    numbersep=5pt,                  
    showspaces=false,                
    showstringspaces=false,
    showtabs=false,                  
    tabsize=2,
    frame=single,                    
    rulecolor=\color{codegray}
}

\lstset{style=miestilo}
% Soporte para tildes y ñ en bloques de código
\lstset{literate={ñ}{{\~n}}1 {á}{{\'a}}1 {é}{{\'e}}1 {í}{{\'i}}1 {ó}{{\'o}}1 {ú}{{\'u}}1}

% Operadores matemáticos personalizados

% Vector en sentido discreto (lista finita de coordenadas): negrita + flecha.
% Uso: \vect{v}, \vect{e}_k, \vect{0}.
\newcommand{\vect}[1]{\vec{\mathbf{#1}}}

% Fecha de versión y comando para capítulos standalone
\providecommand{\verdate}{\the\year.\ifnum\month<10 0\fi\the\month.\ifnum\day<10 0\fi\the\day}
\newcommand{\chapterversion}{%
    \vspace{-1.5cm}\begin{flushright}\small\textit{\color{codegray}Versión~\verdate}\end{flushright}\vspace{0.5cm}%
}

% --- 8. HIPERVÍNCULOS (DEBE SER EL ÚLTIMO PAQUETE) ---
\usepackage[colorlinks=true, linkcolor=utnblue, urlcolor=utnblue, citecolor=utnblue]{hyperref}

% --- 9. REFERENCIAS CRUZADAS ENTRE CAPÍTULOS STANDALONE ---
% Cuando se compila un capítulo standalone, xr-hyper lee los .aux de los
% otros capítulos (si existen) para resolver \ref{} a labels externos.
% En la compilación del libro completo este bloque no se activa.
\ifstandalone
  \usepackage{xr-hyper}
  \makeatletter
  \newcommand{\xrchap}[1]{%
    \edef\xrc@self{\detokenize{Capitulo#1}}%
    \edef\xrc@job{\jobname}%
    \ifx\xrc@self\xrc@job\else
      \IfFileExists{../#1capitulo/Capitulo#1.aux}%
        {\externaldocument{../#1capitulo/Capitulo#1}}{}%
    \fi
  }
  \makeatother
  \xrchap{01}\xrchap{02}\xrchap{03}\xrchap{04}
  \xrchap{05}\xrchap{06}\xrchap{07}\xrchap{08}
  \xrchap{09}\xrchap{10}\xrchap{11}\xrchap{12}
  \xrchap{13}
\fi
% \PassOptionsToPackage{gray}{xcolor}

% fig_choque_sinusoides
% Dos cosenos que comparten muestras bajo Delta t = 1/8 s.
%   - lenta: 1 ciclo/s  -> cos(360*1*x)  (utnblue)
%   - rapida: 9 ciclos/s -> cos(360*9*x) (red!80!black)
% Muestras en x = n/8, n=0..8: y = cos(45*n). Ambas pasan por estos puntos.

\begin{document}
\begin{tikzpicture}[>={Latex[length=3pt,width=3pt]}, font=\scriptsize]
\begin{axis}[
    width=13cm, height=5.6cm,
    axis lines=center,
    axis line style={->, thick},
    % el eje de ordenadas arranca a la altura de las lineas guia, para no
    % atravesar la fila de indices n de abajo
    y axis line style={shorten <=17pt},
    xmin=-0.03, xmax=1.08, ymin=-1.62, ymax=1.45,
    xtick={0,0.125,0.25,0.375,0.5,0.625,0.75,0.875,1},
    xticklabels={$0$,,,,,,,,$1$},
    ytick={-1,0,1},
    every tick label/.style={font=\footnotesize},
    xlabel={$t$ (s)},
    xlabel style={anchor=west, at={(axis description cs:1.0,0.46)}},
    ylabel={},
    enlarge x limits=false,
    clip mode=individual,
    axis on top,
]
    % lineas guia verticales en los instantes de muestreo (sin \foreach dentro del eje)
    \draw[gray!40, thin, dashed] (axis cs:0.125,-1.3) -- (axis cs:0.125,1.3);
    \draw[gray!40, thin, dashed] (axis cs:0.25,-1.3)  -- (axis cs:0.25,1.3);
    \draw[gray!40, thin, dashed] (axis cs:0.375,-1.3) -- (axis cs:0.375,1.3);
    \draw[gray!40, thin, dashed] (axis cs:0.5,-1.3)   -- (axis cs:0.5,1.3);
    \draw[gray!40, thin, dashed] (axis cs:0.625,-1.3) -- (axis cs:0.625,1.3);
    \draw[gray!40, thin, dashed] (axis cs:0.75,-1.3)  -- (axis cs:0.75,1.3);
    \draw[gray!40, thin, dashed] (axis cs:0.875,-1.3) -- (axis cs:0.875,1.3);
    \draw[gray!40, thin, dashed] (axis cs:1,-1.3)      -- (axis cs:1,1.3);

    % sinusoide rapida: 9 ciclos/s
    \addplot[red!80!black, thick, samples=600, domain=0:1] {cos(360*9*x)};
    % sinusoide lenta: 1 ciclo/s
    \addplot[utnblue, thick, samples=200, domain=0:1] {cos(360*1*x)};

    % muestras compartidas
    \addplot[only marks, mark=*, mark size=1.8pt,
             mark options={fill=black, draw=black}]
        coordinates {
            (0,1) (0.125,0.70711) (0.25,0) (0.375,-0.70711) (0.5,-1)
            (0.625,-0.70711) (0.75,0) (0.875,0.70711) (1,1)
        };

    % indices de muestra n = 0..8, alineados con los instantes de muestreo
    \node[anchor=base, font=\footnotesize] at (axis cs:0,-1.42) {$0$};
    \node[anchor=base, font=\footnotesize] at (axis cs:0.125,-1.42) {$1$};
    \node[anchor=base, font=\footnotesize] at (axis cs:0.25,-1.42) {$2$};
    \node[anchor=base, font=\footnotesize] at (axis cs:0.375,-1.42) {$3$};
    \node[anchor=base, font=\footnotesize] at (axis cs:0.5,-1.42) {$4$};
    \node[anchor=base, font=\footnotesize] at (axis cs:0.625,-1.42) {$5$};
    \node[anchor=base, font=\footnotesize] at (axis cs:0.75,-1.42) {$6$};
    \node[anchor=base, font=\footnotesize] at (axis cs:0.875,-1.42) {$7$};
    \node[anchor=base, font=\footnotesize] at (axis cs:1,-1.42) {$8$};
    \node[anchor=west, font=\footnotesize] at (axis cs:1.03,-1.42) {$n$};

    % etiquetas de cada curva
    \node[utnblue, anchor=west, font=\footnotesize]
        at (axis cs:0.265,1.27) {$1$ ciclo/s};
    \node[red!80!black, anchor=west, font=\footnotesize]
        at (axis cs:0.66,1.27) {$9$ ciclos/s};
\end{axis}
\end{tikzpicture}
\end{document}
